\section*{Contexte du Projet de Fin d'Études}

Ce projet de fin d'études s'inscrit dans le cadre d'un stage d'ingénierie au sein de
\textbf{Sopra Steria}, groupe international de conseil et de services numériques.
La mission confiée est de moderniser et de sécuriser le cycle de vie logiciel d'une application
web en adoptant une approche \textbf{DevSecOps} complète, depuis le code source jusqu'à la
production sur le cloud.

La transformation numérique des entreprises impose aujourd'hui des exigences croissantes en
matière de rapidité de livraison, de fiabilité des systèmes, et surtout de \textbf{sécurité
intégrée}. Les organisations qui déploient des applications cloud-native font face à une
complexité croissante~: gestion de conteneurs, orchestration Kubernetes, multiplication des
pipelines CI/CD, et une surface d'attaque qui s'étend à chaque couche de la chaîne de valeur.

\section*{Objectifs du Rapport}

Ce rapport a pour but de~:
\begin{itemize}
  \item Documenter la \textbf{démarche complète} de conception et d'implémentation d'une
        plateforme DevSecOps~;
  \item Démontrer la \textbf{valeur ajoutée} par rapport à l'approche initiale basée sur
        Jenkins sans sécurité~;
  \item Présenter les \textbf{choix techniques} et les justifications architecturales~;
  \item Illustrer les \textbf{résultats concrets} à travers des métriques et artefacts~;
  \item Proposer un \textbf{plan d'amélioration continue} pour les phases suivantes.
\end{itemize}

\section*{Structure du Rapport}

Le rapport est organisé de manière progressive selon quatre chapitres~:

\begin{itemize}
  \item Le \textbf{Chapitre~\ref{chap:contexte}} analyse le contexte, la problématique
        et l'état de l'existant (pipeline Jenkins). Il identifie les lacunes de sécurité
        et justifie la nécessité d'une transformation~;
  \item Le \textbf{Chapitre~\ref{chap:besoins}} présente l'analyse des besoins fonctionnels
        et non fonctionnels du projet, identifie les acteurs, définit les cas d'utilisation,
        et décrit la méthodologie Agile adoptée avec la planification des sprints~;
  \item Le \textbf{Chapitre~\ref{chap:conception}} détaille l'analyse et la conception du
        système~: l'architecture DevSecOps en six couches, l'architecture physique GCP,
        l'architecture logique Kubernetes, les diagrammes de séquence et d'activité,
        ainsi que les choix technologiques~;
  \item Le \textbf{Chapitre~\ref{chap:implementation}} décrit la réalisation du projet
        organisée en quatre sprints Agile~: infrastructure et IaC (Sprint~1), containerisation
        et pipeline CI/CD (Sprint~2), sécurité intégrée et GitOps (Sprint~3), et tests
        automatisés, observabilité et résultats (Sprint~4).
\end{itemize}

Le rapport se conclut par une synthèse des compétences acquises, les limites identifiées
et une roadmap d'amélioration continue à court, moyen et long terme.
